\section{Quickstart: anatomy of a course repo}\label{sec:quickstart}

A course lives in one repository, organised by topic, with this machinery in
a hidden subdirectory. The demo course (\file{course-demo}) is the
reference layout:

\begin{manexsource}
course-demo/
  .course-machinery/        the toolkit (classes, packages, scripts)
  .course-machinery-local/  machinery for THIS course, shared across files
    books.tex               the book registry: every book the course names,
                            with its title and URLs. Loaded into EVERY
                            document, whole
    textbook.tex            the text layer: what a book NAMES things -- the
    openstax.tex            vocabulary, which is exclusive. The first is the
                            course default; a second file is a second book,
                            chosen per document with \usetext
  .MIRRORDIR                optional PDF mirror destination
  derivatives/              a topic (a course has several, same shape)
    problems/               problem bank: one problem per file
    lectures/               LectureNotes masters (+ .latexmkrc, which puts both
                            machinery folders on TEXINPUTS; build-pdfs installs
                            it, one identical copy per build directory)
    facts/                  fact bank: one tagged snippet per file
    images/
  _assessment/              Assessment masters (+ .latexmkrc)
  _syllabus/                CourseDocument masters (+ .latexmkrc)
\end{manexsource}

Three document classes cover every document a course produces -- when you
sit down to write, the only decision is which one:

\begin{itemize}
  \item \file{Assessment} -- anything graded or handed in: tests, finals,
    quizzes, homework (\texttt{[noid]}), worksheets. Clear and standard;
    nothing fancy. Section~\ref{sec:assessments}.
  \item \file{LectureNotes} -- lecture handouts, with a projector view for
    the room. Big type, board space, imported examples.
    Section~\ref{sec:lectures}.
  \item \file{CourseDocument} -- prose documents: syllabus, TA info,
    policies. Coloured headings, TOC, reading measure.
    Section~\ref{sec:prose-documents}.
\end{itemize}

Underneath them, shared packages supply the content machinery -- the same
problem environments, boxes, whitespace commands and fonts everywhere, so
one problem file moves freely between a test and a lecture. You rarely load
any of them yourself; the class does (Part~\ref{part:reference} documents
each).

\subsection{The wiring}

Documents name their class plainly -- \cs{documentclass}\texttt
{\{LectureNotes\}} -- and never spell out where it lives. The
\file{.latexmkrc} in each \file{.tex}-holding directory makes that work: it
walks up to the repo root and puts \file{.course-machinery//} and
\file{.course-machinery-local/} on \env{TEXINPUTS} for that build --- the
two tiers of machinery that are shared across files. Everything else a
document loads sits beside it and is found in the cwd
(section~\ref{sec:extending}). The copies are identical, and you never plant
them yourself: latexmk reads only the rc in its startup directory and never
walks up, so \file{build-pdfs} installs one in each source's directory before
compiling. Creating a new topic folder takes no setup at all.

Cross-file inclusion uses the \file{import} package's relative paths, so one
source serves every consumer. A lecture and the quiz pull the \emph{same}
\file{power-rule-basic.tex} by the path each needs -- \file{../problems/}
beside it, \file{../derivatives/problems/} from the repo root -- and any
images that file loads resolve against the \emph{file's own} directory.

The course-wide content files are two, and the split between them is the
one thing worth learning early. \file{.course-machinery-local/textbook.tex}
is the \emph{text layer}: what the book \emph{names} things --- the result
boxes it uses (APEX's ``Key Idea''), the noun lectures call an example.
That is \emph{exclusive}, since a document cannot use two vocabularies at
once, so \cs{usetext} picks exactly one.
\file{.course-machinery-local/books.tex} is the \emph{book registry}:
each book's title, publisher and URLs. That is not exclusive at all --- a
syllabus names every book the course uses --- so it is loaded into every
document, whole.

Masters \cs{input} neither --- both are found by name and loaded
automatically, so a lecture may sit at any depth and a new master cannot
forget them. Section~\ref{sec:extending} explains the one-question test for
what belongs there, and why the text layer is guarded. Which \emph{section}
of the book a given lecture covers is declared in that lecture, not here
(section~\ref{sec:ref-textref}).

\subsection{The five-minute loop}

Day to day, the toolkit is a three-step loop:

\begin{enumerate}
  \item \textbf{Write or find problems.} One file per problem in a topic's
    \file{problems/}, tagged; \texttt{find-problems -{}-topic ...} to
    rummage (sections \ref{sec:problem-files}, \ref{sec:build-and-find}).
  \item \textbf{Assemble a document.} A master imports problems (with
    points, on a test; as examples, in a lecture) and declares its views
    in the \texttt{\%! views:} comment.
  \item \textbf{Build.} \texttt{build-pdfs} in the document's directory
    compiles everything there -- student PDF plus each declared view,
    routed to the synced mirror if one is configured
    (section~\ref{sec:build-and-find}).
\end{enumerate}

The rest of Part~1 walks these in order; Part~\ref{part:reference} is the
per-file reference; Part~\ref{part:architecture} explains the architecture
-- the load-order contract, the view/layer model, and how to extend the
vocabulary -- for the day you want to change the machinery itself.
